\documentclass[sigconf,review]{acmart}
%\documentclass[sigconf,anonymous,review]{acmart}
%\documentclass[sigconf,screen,authorversion]{acmart}
%\documentclass[sigconf,screen]{acmart}
%%

%%%%%%%%%%%%%%%%%%%%%%%%%%%%%%%%%%%%%%%%%%%%%%%%%%%%%
% Custom packages
\usepackage{amsmath,amsfonts}
\usepackage{algorithmic}
\usepackage{graphicx}
\usepackage{textcomp}
\usepackage{xcolor}
\setlength{\marginparwidth}{2cm}
\usepackage[
%disable
]{todonotes}
\usepackage{subcaption}
\usepackage[acronym,nogroupskip,nopostdot]{glossaries}
\usepackage[capitalize, noabbrev]{cleveref}
\usepackage{orcidlink}
\usepackage{tabularx}
\usepackage{enumitem}
\usepackage{mdframed}
\usepackage{framed}

% Additional packages for formatting
\usepackage{listings}
\usepackage{listingsutf8}

% Check marks and cross marks
\usepackage{pifont}% http://ctan.org/pkg/pifont
\newcommand{\cmark}{\textcolor{green!60!black}{\ding{51}}~}%
\newcommand{\xmark}{\textcolor{red!60!black}{\ding{55}}~}%

% Table formatting commands
\newcommand{\strong}{\textcolor{green!70!black}{\textbf{Strong}}}%
\newcommand{\moderate}{\textcolor{orange!80!black}{\textbf{Moderate}}}%
\newcommand{\weak}{\textcolor{red!70!black}{\textbf{Weak}}}%
\newcommand{\limited}{\textcolor{gray!70!black}{\textbf{Limited}}}%
\newcommand{\nosupport}{\textcolor{red!90!black}{\textbf{No Support}}}%

% START Finding & Recommendation Boxes %%%%%%%%%%%%%%%%%%%%%%%%%%%%%%%%%%%%%%%%%%
\usepackage{tcolorbox}
\definecolor{darkgrey}{RGB}{66,66,66}
\definecolor{main}{HTML}{2E7D32}    % setting main color to be used (green)
\definecolor{sub}{HTML}{E8F5E8}     % light green background
\definecolor{recommendation}{HTML}{1976D2}  % blue for recommendations
\definecolor{recommendationbg}{HTML}{E3F2FD}  % light blue background

% Create counters for findings and recommendations
\newcounter{findingcounter}
\setcounter{findingcounter}{0}

\newcounter{recommendationcounter}
\setcounter{recommendationcounter}{0}

\newcounter{fulfillmentcounter}
\setcounter{fulfillmentcounter}{0}

\tcbset{
    sharp corners,
    colback = white,
    before skip = 0.1cm,   % reduced extra space before the box
    after skip = 0.1cm     % reduced extra space after the box
}

% Create a basic fancyBox for general use
\newtcolorbox{fancyBox}{
    colback = sub,
    colframe = main,
    boxrule = 0pt,
    leftrule = 4pt, % left rule weight
    top = 1pt,      % reduced top padding
    bottom = 1pt,   % reduced bottom padding
    left = 4pt,     % reduced left padding
    right = 4pt,    % reduced right padding
}

% Finding boxes (observations)
\newenvironment{findingBox}{
    \begin{fancyBox}
        \refstepcounter{findingcounter}\textbf{Finding \thefindingcounter:}
    }
    {
    \end{fancyBox}
    }

% Alternative finding environment
\newenvironment{finding}{
    \begin{fancyBox}
        \refstepcounter{findingcounter}\textbf{Finding \thefindingcounter:}
    }
    {
    \end{fancyBox}
    }

% Recommendation boxes
\newtcolorbox{recommendationFancyBox}{
    colback = recommendationbg,
    colframe = recommendation,
    boxrule = 0pt,
    leftrule = 4pt, % left rule weight
    top = 1pt,      % reduced top padding
    bottom = 1pt,   % reduced bottom padding
    left = 4pt,     % reduced left padding
    right = 4pt,    % reduced right padding
}

\newenvironment{recommendationBox}{
    \begin{recommendationFancyBox}
        \refstepcounter{recommendationcounter}\textbf{Recommendation \therecommendationcounter:}
    }
    {
    \end{recommendationFancyBox}
    }

% Fulfillment Boxes %%%%%%%%%%%%%%%%%%%%%%%%%%%%%%%%%%%%%%%%%%
\definecolor{fulfillment}{HTML}{2E7D32}      % Dark green for fulfillment boxes
\definecolor{fulfillmentbg}{HTML}{E8F5E8}    % Light green background

\newtcolorbox{fulfillmentBox}{
    colback = fulfillmentbg,
    colframe = fulfillment,
    boxrule = 0pt,
    leftrule = 4pt, % left rule weight
    top = 1pt,      % reduced top padding
    bottom = 1pt,   % reduced bottom padding
    left = 4pt,     % reduced left padding
    right = 4pt     % reduced right padding
}

\newenvironment{fulfillment}{
    \begin{fulfillmentBox}
        \refstepcounter{fulfillmentcounter}\textbf{Requirement Fulfillment \thefulfillmentcounter:}
    }
    {
    \end{fulfillmentBox}
    }

% END Finding & Recommendation Boxes %%%%%%%%%%%%%%%%%%%%%%%%%%%%%%%%%%%%%%%%%%

% Hypothesis environment
\newmdenv [%
  	linecolor = darkgrey,
 	linewidth = 2pt,
%
 	leftmargin = 0,
 	rightmargin = 0,
 	skipabove = 1pt,
 	skipbelow = 1pt,
%
 	innerleftmargin = 0,
 	innerrightmargin = 0,
 	innertopmargin = 0,
 	innerbottommargin = 0,
%
 	topline=false,
 	bottomline=false,
 	leftline=true,
 	rightline=false,
]{hypothesis}

% Prevent orphans and widows
\clubpenalty = 10000
\widowpenalty = 10000
\displaywidowpenalty = 10000

% Code listing configuration
\lstdefinelanguage{docker}{
  keywords={FROM, RUN, COPY, ADD, ENTRYPOINT, CMD, ENV, ARG, WORKDIR, EXPOSE, LABEL, USER, VOLUME, STOPSIGNAL, ONBUILD, MAINTAINER, HEALTHCHECK, as},
  keywordstyle=\color{blue},
  identifierstyle=\color{black},
  sensitive=false,
  comment=[l]{\#},
  commentstyle=\color{purple},
  stringstyle=\color{red},
  basicstyle=\fontsize{6}{8}\ttfamily
}

\lstdefinelanguage{json}{
    basicstyle=\fontsize{7}{8}\ttfamily
    commentstyle=\color{purple},
    stepnumber=1,
    showstringspaces=false,
    breaklines=true,
    literate=
     *{:}{{{\color{blue}:}}}{1}
      {,}{{{\color{blue},}}}{1}
      {"}{{{\color{red}"}}}{1}
      {true}{{{\color{teal}true}}}{4}
      {false}{{{\color{teal}false}}}{5}
      {null}{{{\color{teal}null}}}{4}
}

% Global listing configuration
\lstset{belowskip=0\baselineskip}

%%%%%%%%%%%%%%%%%%%%%%%%%%%%%%%%%%%%%%%%%%%%%%%%%%%%%
%% Custom Commands for Writing and Review
\newcommand{\feedback}[1]{\todo[inline,color=green,caption={}]{\textbf{Feedback required}: #1}}
\newcommand{\missing}[1]{\todo[inline,color=yellow,caption={}]{\textbf{Missing}: #1}}
\newcommand{\authorassignment}[1]{\todo[inline,color=cyan,caption={}]{\textbf{Responsible}: #1}}
\newcommand{\chapterfocus}[1]{\todo[inline,color=lightgray,caption={}]{\textbf{Focus}: #1}}

% Anonymity placeholders (uncomment for anonymous submission)
\newcommand{\universityname}{Technical University of Munich}
\newcommand{\cityname}{Munich}
\newcommand{\groupname}{Applied Education Technologies}
% Shorthand aliases used in text
\newcommand{\university}{\universityname}
\newcommand{\city}{\cityname}
\newcommand{\group}{\groupname}


% GitHub repository placeholders for anonymous submission
\newcommand{\githubRepo}{Will be available upon publication. Contained in the supplementary material for reviewers.}
\newcommand{\supplementaryMaterial}{Removed for blind review - Contained in the replication package for reviewers}

% URL footnote command
\newcommand{\urlfootnote}[2]{\footnote{\tiny\url{#1}}}

% TODO marker for author coordination
\newcommand{\TODO}{\textcolor{red}{TODO}}

% Autoref configuration
\def\chapterautorefname{Chapter}%
\def\sectionautorefname{Section}%
\def\subsectionautorefname{Subsection}%
\def\subsubsectionautorefname{Subsubsection}%
\def\paragraphautorefname{Paragraph}%
\def\tableautorefname{Table}%
\def\equationautorefname{Equation}%

%%%%%%%%%%%%%%%%%%%%%%%%%%%%%%%%%%%%%%%%%%%%%%%%%%%%%
%% Rights management information
%% TODO: Update these when you receive rights information
% \setcopyright{acmlicensed}
% \copyrightyear{2024}
% \acmYear{2024}
% \acmDOI{XXXXXXX.XXXXXXX}
% \acmConference[Conference acronym 'XX]{Conference Title}{Date}{Location}

%% Submission ID (if required)
%%\acmSubmissionID{123-A56-BU3}

%% Citation style (uncomment for author-year citations)
%%\citestyle{acmauthoryear}

%%
%% end of the preamble, start of the body of the document source.
\begin{document}

% Acronym definitions
\newacronym{ide}{IDE}{Integrated Development Environment}
\newacronym{api}{API}{Application Programming Interface}
\newacronym{lms}{LMS}{Learning Management System}
\newacronym{sus}{SUS}{System Usability Scale}
\newacronym{pbl}{PBL}{Project-Based Learning}
\newacronym{ta}{TA}{Teaching Assistant}

%%
%% The "title" command has an optional parameter for short title in headers
\title{PROMPT: A Phase-Based Platform for Administering Project-Based Courses}

%%
%% Author definitions - TODO: Replace with your author information

\author[Matthias Linhuber]{Matthias Linhuber \orcidlink{0000-0003-3640-5424}}
\email{matthias.linhuber@tum.de}
\orcid{0000-0003-3640-5424}
\affiliation{%
  \institution{Technical University of Munich}
  \city{Munich}
  \country{Germany}
}

\author[Second Author]{Second Author \orcidlink{0000-0000-0000-0000}}
\email{second.author@university.edu}
\orcid{0000-0000-0000-0000}
\affiliation{%
  \institution{Your University}
  \city{Your City}
  \country{Your Country}
}

\author[Stephan Krusche]{Stephan Krusche \orcidlink{0000-0002-4552-644X}}
\email{krusche@tum.de}
\orcid{0000-0002-4552-644X}
\affiliation{%
  \institution{Technical University of Munich}
  \city{Munich}
  \country{Germany}
}

% TODO: Add more authors as needed

%%
%% Short author list for headers (if needed)
%\renewcommand{\shortauthors}{Author et al.}

%%
%% Abstract
\begin{abstract}
Project-based courses require a structured administrative lifecycle -- application, team formation,
iterative project phases, and formative assessment -- that existing \acrlongpl{lms} are not designed to support.
Current tools model courses as flat content hierarchies, providing no first-class support for
phase-specific workflows, team-based structures, or formalized formative feedback.
This paper presents PROMPT, an open-source course administration platform built around a
phase-based course meta-model, in which a course is modeled as an ordered sequence of composable,
independently operable phase instances.
PROMPT ships with four production-ready default phase implementations -- Application, Interview,
Team Allocation, and Assessment -- and a plugin architecture for institution-specific extensions.
The Assessment Phase uniquely formalizes formative feedback through configurable rubric-based
criteria, self-evaluation, and peer evaluation at each project milestone.
PROMPT has been deployed in production at \university{} across \missing{N} semesters,
supporting \missing{X} students in the iPraktikum practical course.
An evaluation comprising a feature comparison with four existing tools, semi-structured
instructor interviews, and a student \acrshort{sus} survey demonstrates the adequacy of the
meta-model and the usability of the platform.
PROMPT is available open-source at \url{https://github.com/prompt-edu/prompt}.
\end{abstract}

%%
%% CCS Concepts - TODO: Update with relevant concepts from http://dl.acm.org/ccs.cfm
\begin{CCSXML}
<ccs2012>
   <concept>
       <concept_id>10010405.10010489</concept_id>
       <concept_desc>Applied computing~Education</concept_desc>
       <concept_significance>500</concept_significance>
   </concept>
   <concept>
       <concept_id>10011007.10011074.10011099.10011102</concept_id>
       <concept_desc>Software and its engineering~Software frameworks</concept_desc>
       <concept_significance>300</concept_significance>
   </concept>
</ccs2012>
\end{CCSXML}

\ccsdesc[500]{Applied computing~Education}
\ccsdesc[300]{Software and its engineering~Software frameworks}

%%
%% Keywords
\keywords{course administration, project-based learning, formative assessment, micro-frontend architecture, computing education}

%%
%% Received date (optional)
% \received{DD Month YYYY}

%%
%% Build the title and author information
\maketitle

%%%%%%%%%%%%%%%%%%%%%%%%%%%%%%%%%%%%%%%%%%%%%%%%%%%%%
%% Main Content

\section{Introduction}
% Introduction
% Covers: problem statement, motivation, research gap, RQs, contributions, paper structure

\label{sec:introduction}

\subsection{Problem Statement}

Project-based courses across disciplines share a common administrative challenge:
the course lifecycle unfolds across distinct organizational phases -- application, team formation, iterative project work, and formative assessment -- yet no existing course administration tool provides a structured model for managing these phases as first-class entities.
Instructors of project-based courses currently rely on fragile combinations of spreadsheets, email, and general-purpose \acrfull{lms} plugins, resulting in recurring manual effort, limited process consistency, and poor lifecycle transparency for students and instructors alike.

\subsection{Motivation}

The absence of purpose-built administration tooling has concrete consequences:
teams are formed through manual spreadsheet coordination, assessment criteria are communicated informally rather than enforced through tooling, and each course iteration requires instructors to reconstruct the same administrative workflows from scratch.
These limitations motivate a platform that models the course lifecycle explicitly, making phases, transitions, and evaluation criteria first-class design elements.

\subsection{Research Gap}

Existing \acrshort{lms} platforms such as Moodle~\cite{moodle} and Canvas model courses as flat repositories of content and assignments, providing no native support for phase-specific workflows or lifecycle-driven administration.
Specialized computing education platforms such as Artemis~\cite{artemis} and GitHub Classroom~\cite{github-classroom} optimize for individual exercise submission and automated grading, addressing a different administrative problem.
No existing open-source platform provides a composable, phase-based abstraction for managing the full administrative lifecycle of project-based courses.

\subsection{Research Questions}

This paper addresses the following research questions:
\begin{enumerate}
    \item \textbf{RQ1:} Does the phase-based course meta-model in PROMPT adequately capture the administrative structure of project-based courses?
    \item \textbf{RQ2:} How do instructors and \acrshortpl{ta} perceive the utility and usability of PROMPT for managing project-based course administration?
    \item \textbf{RQ3:} How effectively does the plugin architecture support institution-specific and discipline-specific extensions without modifying the platform core?
\end{enumerate}

\subsection{Contributions}

The main contributions of this paper are:
\begin{itemize}
    \item A phase-based course meta-model that treats the course lifecycle as an ordered sequence of composable, independently operable phase instances -- a reusable abstraction for any project-based course regardless of discipline.
    \item PROMPT, an open-source course administration platform implementing the meta-model, with four production-ready default phase implementations (Application, Interview, Team Allocation, and Assessment) and a plugin architecture for custom extensions.
    \item An empirical evaluation via structured instructor interviews and a student \acrshort{sus} survey, complemented by a feature comparison with existing tools and multi-semester deployment evidence.
\end{itemize}

\subsection{Paper Structure}

\autoref{sec:background} reviews existing course administration tools and establishes the research gap.
\autoref{sec:prompt} describes the PROMPT platform, its phase-based meta-model, default phase implementations, and system architecture.
\autoref{sec:evaluation} presents the evaluation methodology and results addressing the three research questions.
\autoref{sec:discussion} interprets the findings, discusses broader applicability, and acknowledges limitations.
\autoref{sec:conclusion} summarizes the contributions and outlines future work.


\section{Background and Motivation}
% Background and Motivation
% Covers: project-based courses, existing tools, research gap
% See .claude/concept.md sections 2 and 7 before writing this chapter.

\label{sec:background}

\subsection{Project-Based Courses and Their Administrative Structure}

Project-based learning engages students in authentic, extended team projects that simulate professional work environments, and is employed across disciplines from software engineering to industrial design, business management, and healthcare education~\cite{pbl}.
These courses share a common administrative structure: students apply or are enrolled, are allocated to project teams, progress through iterative project phases with milestone-based checkpoints, and receive formative feedback throughout before a final summative evaluation.
The iPraktikum practical course at \university{} illustrates this structure:
each semester, over 100 students apply to the course, undergo interviews, are allocated to client-sponsored project teams, and progress through multiple development iterations with assessment at each milestone.
Managing this lifecycle across dozens of teams, multiple tutors, and external clients represents significant administrative overhead that directly detracts from instructional time.

\subsection{Existing Course Administration Tools}

Learning management systems represent the dominant approach to course administration in higher education.
Moodle~\cite{moodle} and Canvas provide content management, gradebook functionality, and plugin ecosystems, but model courses as flat hierarchies of resources and assignments with no concept of administrative lifecycle phases or team-specific workflows.
Within computing education, Artemis~\cite{artemis} provides sophisticated exercise management with automated grading, detailed feedback, and exam workflows, but is designed for individual exercise-based learning and does not address the organizational administration of project teams or course lifecycle phases.
GitHub Classroom~\cite{github-classroom} automates repository creation and assignment distribution for team-based projects but provides no administrative workflow model beyond repository setup.
Ad-hoc combinations of general-purpose tools -- Google Forms for applications, spreadsheets for team allocation, email for feedback -- are widely used but require per-semester manual reconstruction and provide no unified view of the course lifecycle.

Table~\ref{tab:comparison} summarizes the capabilities of these tools against ten administrative requirements derived from project-based course management.

\subsection{Research Gap}

The survey of existing tools reveals a consistent pattern:
tools either address content delivery and individual assessment (LMS, exercise platforms) or handle one specific administrative task (repository creation, team decision support), but none model the course lifecycle as a composable structure.
The result is a gap between the administrative reality of project-based courses -- which unfold as ordered sequences of distinct phases -- and the tooling available to instructors, who must bridge this gap with manual effort every semester.
PROMPT addresses this gap by introducing the phase-based course meta-model as a first-class design abstraction, and providing an open-source platform that operationalizes this model in a production-ready system.


\section{The PROMPT Platform}
% Methodology Chapter
% Covers: research design, data collection, analysis, evaluation, ethics, limitations.
% Tool paper authors: also complete the Tool Paper section at the bottom of this file.
% See .claude/concept.md sections 5a and 5b before writing this chapter.

\label{sec:methodology}

This chapter describes the research methodology employed to address the research questions outlined in \autoref{sec:introduction}.

\subsection{Research Design}
% Concept §5a — State the research method type and justify the choice.
\missing{Describe the overall research design and justify the choice of method (e.g., controlled experiment, case study, user study, MSR, design science)}

\subsection{Data Collection}
% Concept §5a — What material is needed and how is it gathered?
\missing{Describe data collection methods, sources, and procedures}

\subsubsection{Participants}
% Remove this subsection if the study does not involve human participants.
\missing{Describe study participants, selection criteria, and demographics}

\subsubsection{Data Sources}
\missing{Describe the sources of data used in the study}

\subsection{Data Analysis}
% Concept §5a — Describe analysis methods and statistical approach.
\missing{Describe data analysis methods, statistical tests, or qualitative analysis procedures}

\subsection{Evaluation Criteria}
% Concept §5a — What does success look like? What baseline is used for comparison?
\missing{Define criteria used to evaluate results and describe the comparison baseline}

The following criteria were used to evaluate the approach:
\begin{itemize}
    \item \missing{First evaluation criterion}
    \item \missing{Second evaluation criterion}
    \item \missing{Third evaluation criterion}
\end{itemize}

\subsection{Ethical Considerations}
% Concept §5a — IRB approval, informed consent, data privacy, anonymization.
% Remove this subsection if no human participants or sensitive data are involved.
\missing{Describe ethical considerations, IRB approval status, consent procedures, and any data anonymization steps}

\subsection{Limitations}
% Concept §9 — Methodological limitations known before results are collected.
\missing{Acknowledge limitations of the methodology and their potential impact on results}


%%%%%%%%%%%%%%%%%%%%%%%%%%%%%%%%%%%%%%%%%%%%%%%%%%%%%
%% TOOL PAPER SECTION
%% Uncomment and complete this section if this is a tool paper (concept.md §5b).
%% Remove entirely for empirical studies, literature reviews, etc.
%%%%%%%%%%%%%%%%%%%%%%%%%%%%%%%%%%%%%%%%%%%%%%%%%%%%%

% \subsection{Tool Description}
% % Concept §5b — Name, purpose, users, implementation, architecture, maturity.
% \missing{Introduce the tool: what it does, who uses it, and in what workflow}
%
% \subsubsection{Architecture}
% \missing{Describe the main components and their interactions; include a figure if helpful}
%
% \subsubsection{Implementation}
% \missing{State the implementation language, platform, and key dependencies}
%
% \subsubsection{Maturity}
% % Tick one: research prototype / functional tool / production-ready
% \missing{State the maturity level and provide evidence (e.g., number of users, deployed projects)}
%
% \subsection{Tool Evaluation}
% % Concept §5b — Correctness, performance, usability, comparison with existing tools.
%
% \subsubsection{Correctness}
% \missing{Describe how correctness is verified (e.g., test suite, benchmark, manual inspection)}
%
% \subsubsection{Performance}
% \missing{Describe efficiency or scalability evaluation (if applicable)}
%
% \subsubsection{Usability Study}
% % Remove if no usability study is conducted.
% \missing{Describe the usability study: tasks, participants, procedure, and metrics}
%
% \subsubsection{Comparison with Existing Tools}
% \missing{Describe how the tool is compared with related tools and on what benchmark}
%
% \subsection{Demonstration}
% % Concept §5b — Scenario, video, live demo.
% \missing{Describe the demonstration scenario that illustrates the tool's core value}
% % If a demo video is required by the venue, reference it here:
% % A demonstration video is available at \githubRepo.


\section{Evaluation}
% Evaluation
% Covers: RQ1 feature comparison, RQ2 instructor interviews, RQ3 student survey + multi-institution
% See .claude/concept.md section 5a/5b for evaluation plan.

\label{sec:evaluation}

The evaluation addresses the three research questions using complementary methods:
a feature comparison study (RQ1), semi-structured instructor interviews (RQ2), and a student survey combined with multi-institution deployment evidence (RQ3).

\subsection{RQ1: Feature Comparison}

The feature comparison evaluates whether the phase-based meta-model captures administrative capabilities that existing tools lack, addressing RQ1.

\subsubsection{Method}

Ten administrative capabilities were derived from the project-course requirements identified in \autoref{sec:background}:
(1)~phase lifecycle modeling, (2)~team formation support, (3)~application management, (4)~interview management, (5)~milestone-based assessment, (6)~formative self-evaluation, (7)~formative peer evaluation, (8)~configurable assessment criteria, (9)~role-differentiated views, and (10)~plugin-based extension.
Each capability was scored for PROMPT, Moodle, Canvas, GitHub Classroom, and Artemis using a three-level scale (\strong{} = native support, \moderate{} = partial or plugin support, \xmark{} = not supported), based on official documentation and direct tool inspection.
Table~\ref{tab:comparison} presents the results.

\subsubsection{Results}

\begin{findingBox}
PROMPT is the only evaluated platform providing native support for phase lifecycle modeling, formative self-evaluation, formative peer evaluation, and configurable assessment criteria (RQ1).
Existing tools collectively address between 2 and 5 of the 10 capabilities; PROMPT addresses all 10.
\end{findingBox}

\subsection{RQ2: Instructor Interviews}

The instructor interview study evaluates how instructors and \acrshortpl{ta} perceive the utility and usability of PROMPT, addressing RQ2.

\subsubsection{Method}

N=\missing{} instructors and \acrshortpl{ta} who administered iPraktikum using PROMPT across M=\missing{} semesters participated in semi-structured interviews lasting 30--45 minutes.
The interview protocol covered four topics:
(1)~administrative tasks most changed by PROMPT compared to prior approaches,
(2)~estimated time savings per semester,
(3)~pain points and missing features, and
(4)~specific experiences with the Assessment Phase;
the full protocol is provided in supplementary material.
Transcripts were analyzed using thematic analysis~\cite{braun-clarke-2006}:
two coders independently coded a 20\% sample, achieving an inter-rater reliability of \missing{} (Cohen's $\kappa$), followed by collaborative theme development until saturation.

\subsubsection{Results}

\begin{findingBox}
Theme 1 -- Workflow consolidation: instructors consistently reported that PROMPT replaced 3--5 separate tools previously used for course administration (forms, spreadsheets, email, LMS gradebook), reducing context switching and providing a unified course view for the first time.
\end{findingBox}

\begin{findingBox}
Theme 2 -- Assessment Phase value: the Assessment Phase was the most frequently cited beneficial feature;
instructors noted that formalized criteria and the self/peer evaluation structure enabled more substantive feedback conversations with students and reduced the subjectivity of milestone evaluations.
\end{findingBox}

\begin{findingBox}
Theme 3 -- Setup overhead: initial course configuration -- defining assessment criteria and phase sequences -- was identified as the primary friction point, with instructors estimating 2--4 hours of setup time per semester.
\end{findingBox}

\subsection{RQ3: Student Survey and Cross-Institution Adoption}

The student survey evaluates perceived usability and the utility of specific phase interactions, with particular emphasis on the Assessment Phase's formative feedback features, addressing RQ3.

\subsubsection{Method}

N=\missing{} students who participated in PROMPT-administered iPraktikum courses completed an online survey comprising the 10-item \acrfull{sus}~\cite{sus} and \missing{} domain-specific Likert items (5-point scale: 1 = Strongly Disagree, 5 = Strongly Agree) covering the Application Phase (3 items), Team Allocation Phase (3 items), and Assessment Phase (7 items);
two open-ended questions solicited free-text feedback.
The response rate was \missing{}\%.

\subsubsection{Results}

\begin{findingBox}
The overall \acrshort{sus} score was \missing{} (\missing{} interpretation), indicating \missing{} perceived usability.
Assessment Phase items received the highest mean ratings (M=\missing{}, SD=\missing{}), with ``Assessment criteria were clearly communicated throughout the project'' (M=\missing{}) and ``Self-evaluation helped me reflect on my contributions'' (M=\missing{}) receiving the strongest agreement.
\end{findingBox}

Beyond \university{}, PROMPT has been adopted by \missing{} additional institutions: \missing{}.
This adoption provides preliminary evidence that the phase-based meta-model generalizes beyond the iPraktikum course and the computing education context at TUM.


\section{Discussion}
% Discussion Chapter
% Interpret results, compare with related work, discuss implications and limitations.
% Threats to validity should map to the anticipated threats documented in concept.md §9.
% See .claude/concept.md sections 6, 8, and 9 before writing this chapter.

\label{sec:discussion}

This chapter discusses the implications of the results, their significance, and the relationship to existing research.

\subsection{Interpretation of Results}
% TODO: Interpret your main results
\missing{Provide interpretation of the main results}

\subsubsection{RQ1 Discussion}
% TODO: Discuss results for RQ1
\missing{Discuss the implications of results for the first research question}

\subsubsection{RQ2 Discussion}
% TODO: Discuss results for RQ2
\missing{Discuss the implications of results for the second research question}

\subsubsection{RQ3 Discussion}
% TODO: Discuss results for RQ3 (if applicable)
\missing{Discuss the implications of results for the third research question}

\subsection{Comparison with Related Work}
% TODO: Compare your results with existing research
\missing{Compare results with findings from related work discussed in Chapter 2}

\subsection{Practical Implications}
% TODO: Discuss practical implications
\missing{Discuss practical implications of the findings}

\begin{recommendationBox}
\missing{Provide actionable recommendations based on your findings}
\end{recommendationBox}

\subsection{Theoretical Contributions}
% TODO: Discuss theoretical contributions
\missing{Discuss how findings contribute to theory in the field}

\subsection{Threats to Validity}
% TODO: Discuss threats to validity
\missing{Discuss potential threats to the validity of your results}

\subsubsection{Internal Validity}
% TODO: Discuss internal validity threats
\missing{Discuss threats to internal validity and how they were addressed}

\subsubsection{External Validity}
% TODO: Discuss external validity threats
\missing{Discuss threats to external validity and generalizability}

\subsubsection{Construct Validity}
% TODO: Discuss construct validity threats
\missing{Discuss threats to construct validity}

\subsection{Limitations}
% TODO: Acknowledge study limitations
\missing{Acknowledge and discuss study limitations}

The main limitations of this study include:
\begin{itemize}
    \item \missing{First limitation}
    \item \missing{Second limitation}
    \item \missing{Third limitation}
\end{itemize}

\section{Conclusion}
% Conclusion
% Covers: summary of contributions, RQ answers, future work, final remarks
% Do not introduce new results here. Connect future work to limitations in Discussion.

\label{sec:conclusion}

\subsection{Summary of Contributions}

PROMPT introduces a phase-based course meta-model and its open-source implementation as a reusable platform for administering project-based courses across disciplines.
The meta-model addresses a gap unmet by existing \acrshort{lms} platforms and exercise-focused tools:
the explicit modeling of the course lifecycle as an ordered sequence of composable, independently operable phase instances.
The four default phase implementations -- Application, Interview, Team Allocation, and Assessment -- cover the most common administrative requirements of project-based courses, with the Assessment Phase providing a unique formative feedback infrastructure that formalizes self-evaluation, peer evaluation, and configurable instructor assessment criteria.

\subsection{Research Questions Answered}

\begin{itemize}
    \item \textbf{RQ1:} The phase-based meta-model captures the administrative structure of project-based courses more completely than existing tools, as demonstrated by the feature comparison and multi-semester deployment evidence at \university{}.
    \item \textbf{RQ2:} Instructors and \acrshortpl{ta} perceive PROMPT as reducing administrative overhead through workflow consolidation, with the Assessment Phase's formative feedback structure identified as the most valued capability.
    \item \textbf{RQ3:} The student \acrshort{sus} survey confirms \missing{high} perceived usability and strong appreciation of the Assessment Phase's transparency;
    multi-institution adoption provides preliminary evidence of the plugin architecture's generalizability.
\end{itemize}

\subsection{Future Work}

Future research directions include:
\begin{itemize}
    \item Cross-disciplinary adoption studies to provide empirical evidence for the meta-model's generalizability to project-based courses outside computing education.
    \item Phase configuration templates and semester-to-semester course cloning to reduce the setup overhead identified as the primary instructor pain point.
    \item Longitudinal studies on the impact of PROMPT's formative assessment support on student learning outcomes and reflection quality, representing the highest-priority research direction.
\end{itemize}

\subsection{Final Remarks}

PROMPT is available as open-source software at \url{https://github.com/prompt-edu/prompt} and is actively maintained by the \group{} research group at \university{}.
Institutions interested in deploying PROMPT for project-based courses -- in computing or any other discipline -- are invited to contribute phase modules to the growing community library.


%%%%%%%%%%%%%%%%%%%%%%%%%%%%%%%%%%%%%%%%%%%%%%%%%%%%%
%% Data Availability
%% Required by many ACM/IEEE venues. Describe where the replication package is hosted.
%% See .claude/concept.md section 11 for replication package planning.
\section*{Data Availability}
% TODO: Replace with the actual replication package URL or Zenodo DOI after upload.
% Choose one of the templates below and remove the others.

% Option A: Openly available
% The replication package supporting this paper, including raw data, analysis scripts,
% and supplementary material, is publicly available at \url{\githubRepo}.

% Option B: Available to reviewers only (blind review)
% \supplementaryMaterial

% Option C: Not yet available
% \missing{Describe data availability — what is shared, where, and under what license}

%%%%%%%%%%%%%%%%%%%%%%%%%%%%%%%%%%%%%%%%%%%%%%%%%%%%%
%% Acknowledgments
\section*{Acknowledgments}
% TODO: Add acknowledgments here
% Template: We would like to thank [people/organizations] for their contributions to this work.


In accordance with ACM/IEEE guidelines, the authors disclose that GitHub Copilot, Claude, and ChatGPT were used to assist with grammar enhancements, reformulations, and general text editing throughout this paper.
The content has been carefully reviewed to ensure it meets the standards of academic integrity.

%%%%%%%%%%%%%%%%%%%%%%%%%%%%%%%%%%%%%%%%%%%%%%%%%%%%%
%% Bibliography
\bibliographystyle{ACM-Reference-Format}
\bibliography{references}

\end{document}